\emph{Preliminaries.} The blades are a basis of $\Cl_{\Q}(2,2)$ (Theorem \claimref{synthesis.signed-directions-and-the-integer-clifford-algebra}); the adjoint $\dagger$ is the linear extension of the inverse on blades, an anti-automorphism with $\langle x^{\dagger}\rangle_{0} = \langle x\rangle_{0}$. By Proposition \claimref{tower.ladder-pairs-satisfy-car-vectors}, $Q(v_{j}) = Q(w_{j}) = 0$, $v_{j}w_{j} + w_{j}v_{j} = 1$, ladder elements of different modes anticommute, and $w_{j} = v_{j}^{\dagger}$. The relations $s_{i}^{\dagger}s_{j} = \delta_{ij}P_{0}$, $\sum_{i}s_{i}s_{i}^{\dagger} = 1$, $\langle P_{0}\rangle_{0} = \tfrac14$ and $P_{0}xP_{0} \in \Q P_{0}$ for every blade $x$ follow from these and are verified by the check.

\emph{(i).} Define $m(x)_{ij}$ by $s_{i}^{\dagger}xs_{j} = m(x)_{ij}P_{0}$ (the left side is $P_{0}(s_{i}^{\dagger}xs_{j})P_{0}$, a multiple of $P_{0}$). Then $s_{i}^{\dagger}xys_{j} = \sum_{k}(s_{i}^{\dagger}xs_{k})(s_{k}^{\dagger}ys_{j}) = \sum_{k}m(x)_{ik}m(y)_{kj}P_{0}$, so $m$ is multiplicative; $x = \sum_{i,j}s_{i}s_{i}^{\dagger}xs_{j}s_{j}^{\dagger} = \sum_{i,j}m(x)_{ij}\,s_{i}s_{j}^{\dagger}$ (as $s_{i}P_{0} = s_{i}$), so $m$ is injective, and onto by dimension $16$; $s_{i}^{\dagger}x^{\dagger}s_{j} = (s_{j}^{\dagger}xs_{i})^{\dagger} = m(x)_{ji}P_{0}$, so $m(x^{\dagger})$ is the transpose. With $y_{r} = \sum_{i}r_{i}s_{i}$, $\Pi_{r} = \sum_{i,j}r_{i}r_{j}s_{i}s_{j}^{\dagger}/(r \cdot r)$ has matrix $rr^{T}/(r \cdot r)$. A symmetric idempotent rational matrix is the orthogonal projection onto its range, a rational subspace; that subspace has a rational basis of pairwise orthogonal vectors $r_{a}$ (orthogonalise without normalising), and the projection is $\sum_{a}r_{a}r_{a}^{T}/(r_{a} \cdot r_{a})$. So every projector is a sum of pairwise orthogonal atoms $\Pi_{r_{a}}$, and those of rank one are the $\Pi_{r}$. If $\mu$ is additive, induction on the number of summands gives $\mu(E) = \sum_{a}\mu(\Pi_{r_{a}})$, in particular $\sum_{a=1}^{4}f(r_{a}) = \mu(1)$ on a frame. Conversely a function $f$ on rays with constant frame sums defines $\mu(E) = \sum_{a}f(r_{a})$: two orthogonal bases of the range of $E$, each completed by one fixed orthogonal basis of the complement, give two frames, so the sum does not depend on the basis, and additivity is then immediate.

\emph{(ii).} By \claimref{std-algebra.quaternions} the elements $1, B_{1}, B_{2}, B_{3}$ span the even subalgebra with $B_{k}^{2} = -1$, $B_{1}B_{2} = B_{3}$ and cyclically. For $x = x_{r}$, $y = x_{r'}$: $x\tilde{x} = r \cdot r$ and $\langle x\tilde{y}\rangle_{0} = r \cdot r'$ (the check compares the component formula with the product of the algebra). So $xy^{-1} = x\tilde{y}/(r' \cdot r')$ has no scalar part exactly when $r \cdot r' = 0$. Let $x_{1}, \dots, x_{4}$ be a frame and $a = x_{2}x_{1}^{-1}$, $b = x_{3}x_{1}^{-1}$, $c = x_{4}x_{1}^{-1}$, bivectors. Then $ab^{-1} = x_{2}x_{3}^{-1}$ is a bivector, so $\langle a\tilde{b}\rangle_{0} = 0$: $a$ and $b$ are orthogonal as elements of the three-dimensional space of bivectors, they anticommute, and $ab$ is a bivector orthogonal to both; so is $c$, hence $c$ is a multiple of $ab$. Conversely for orthogonal bivectors $a, b$ the four elements $x, ax, bx, abx$ have pairwise quotients $a, b, ab$ and their products, all bivectors up to scale.

\emph{(iii).} Subtracting the constant $\tfrac14\mu(1)$ (which has frame sums $\mu(1)$ and is the quadratic form of a multiple of the identity) we may assume frame sums $0$. Let $G$ be the group of nonzero elements of $\Cl_{\Q}^{+}(3,0)$ modulo $\Q^{\times}$, $[x]$ the class of $x$ in the group algebra $\Q[G]$, and for orthogonal nonzero rational bivectors $a, b$ put $\kappa_{a,b} = 1 + [a] + [b] + [ab]$. Let $J$ be the span of the elements $\kappa_{a,b}[x]$. Since $[y]\kappa_{a,b} = \kappa_{a',b'}[y]$ with $a' = yay^{-1}$, $b' = yby^{-1}$ (orthogonal bivectors), $J$ is a two-sided ideal; put $A = \Q[G]/J$. By (ii) a function with frame sums $0$ is a linear map on $\Q[G]$ that vanishes on $J$, that is, a linear map $A \to V$. Let $R(x)$ be the matrix of $n \mapsto xnx^{-1}$ on the bivectors in the basis $B_{k}$, a rotation with rational entries; $R$ extends to an algebra map $\Q[G] \to \Mat_{3}(\Q)$. For a bivector $a$, $R(a) = 2P_{a} - 1$ with $P_{a}$ the orthogonal projection on the line of $a$ (a half turn), and for an orthogonal triple $P_{a} + P_{b} + P_{ab} = 1$, so $R(\kappa_{a,b}) = 0$ and $R$ is defined on $A$.

\emph{Step 1.} The $24$ rays of $O$ are closed under the product and inverse (check). The nine bivectors among them form four orthogonal triples, and the $24$ sets $\{x, ax, bx, abx\}$, $x \in O$, with $a, b$ in one of these triples, are $24$ relations among the classes $[x]$, $x \in O$, in $A$; their rank is $15$ (check). So the span $W$ of the $[x]$, $x \in O$, has dimension at most $9$. It is a subalgebra, and $R(O)$ spans $\Mat_{3}(\Q)$ (check: rank $9$). Hence $R$ restricts to an algebra isomorphism $W \to \Mat_{3}(\Q)$, and $W$ contains the unit $[1]$ of $A$. We identify $W$ with $\Mat_{3}(\Q)$ by $R$. Let $G_{W} = \{x \in G : [x] \in W\}$; for $x \in G_{W}$, $[x] = R(x)$. It is a subgroup: it is closed under products, and if $[x] \in W$ then left multiplication by $[x]$, injective on the finite-dimensional $W$ because $[x]$ is invertible in $A$, is onto $W$, so $[x]z = 1$ for some $z \in W$ and $z = [x^{-1}]$.

\emph{Step 2.} The $34$ frames of the check, with the rational coefficients listed there, satisfy
\[
\sum_{F}q_{F}\sum_{y \in F}[y] = [2 + B_{1}] + \sum_{s \in O}t_{s}[s]
\]
in $\Q[G]$ (an identity between formal combinations of rays, verified exactly; $18$ of the $t_{s}$ are nonzero and $\sum_{F}q_{F} = \tfrac25$). Each inner sum is an element $\kappa_{a,b}[x]$ of $J$, so $[2 + B_{1}] = -\sum_{s}t_{s}[s] \in W$. Conjugating by the elements of $O$ gives the six rotors $2 \pm B_{k}$ (check), so they lie in $G_{W}$.

\emph{Step 3.} Let $y$ have integer components and $5 \mid y\tilde{y}$. Then $\tilde{\pi}y \in 5\Z^{4}$ for one of the six $\pi = 2 \pm B_{k}$: the condition depends only on $y$ modulo $5$ and holds for each of the $145$ classes of norm $0$ modulo $5$ (check). Then $y = \pi y'$ with $y' = \tilde{\pi}y/5$ of integer components and $y'\tilde{y}' = y\tilde{y}/5$. By induction every $y$ with integer components and $y\tilde{y} = 5^{t}$ is a product of $t$ of the six rotors and an element of norm $1$, which is $\pm 1$ or $\pm B_{k}$, in $O$. So every such $y$ lies in $G_{W}$.

\emph{Step 4.} Let $n$ be a nonzero bivector with integer components, $N = n\tilde{n}$, so $n^{-1} = -n/N$, and write $N = 5^{e}N'$ with $5 \nmid N'$. In the finite commutative ring $\Z[B_{1}]/(N')$ the element $2 + B_{1}$ is invertible (its product with $2 - B_{1}$ is $5$), so $(2 + B_{1})^{s} \equiv 1$ for some $s \ge 1$: $g = (2 + B_{1})^{s} = \alpha + \beta B_{1}$ with integers $\alpha, \beta$, $N' \mid \beta$, $\alpha^{2} + \beta^{2} = 5^{s}$. Modulo $5$, $(2 + B_{1})^{2} = 3 + 4B_{1} \equiv 4(2 + B_{1})$, so $g \equiv 4^{s-1}(2 + B_{1})$ and neither $\alpha$ nor $\beta$ is divisible by $5$; in particular $\alpha\beta \ne 0$, and $R(g)$ is the rotation about the first axis by an angle $\varphi$ with $\sin\varphi = 2\alpha\beta/5^{s} \ne 0$. Now $ngn^{-1} = \alpha - (\beta/N)\,nB_{1}n$, and $nB_{1}n$ has integer components, so $y = 5^{e}\alpha - (\beta/N')\,nB_{1}n$ has integer components, norm $5^{2e+s}$, and the same class as $ngn^{-1}$. By Step 3, $g' = ngn^{-1} \in G_{W}$, and $g \in G_{W}$ by Step 2. In $A$, $[n][g] = [g'][n]$, that is $[n]R(g) = R(n)R(g)R(n)^{-1}[n]$, and $R(n)^{2} = 1$. So $v_{n} = R(n)[n]$ satisfies $v_{n}R(g) = R(g)v_{n}$. The same holds with $B_{2}$ in place of $B_{1}$, with a rotation about the second axis by an angle with nonzero sine. With matrix units $E_{ij}$ of $W$: $E_{11} = (R(g)^{2} - 2\cos\varphi\,R(g) + 1)/(2 - 2\cos\varphi)$ and $E_{32} - E_{23} = (R(g) - E_{11} - \cos\varphi\,(1 - E_{11}))/\sin\varphi$, so $v_{n}$ commutes with $E_{11}$, $E_{32} - E_{23}$ and likewise with $E_{22}$, $E_{13} - E_{31}$. These generate $\Mat_{3}(\Q)$: $E_{33} = 1 - E_{11} - E_{22}$, $E_{32} = (E_{32} - E_{23})E_{22}$, $E_{23} = -(E_{32} - E_{23})E_{33}$, $E_{13} = (E_{13} - E_{31})E_{33}$, $E_{31} = -(E_{13} - E_{31})E_{11}$, $E_{12} = E_{13}E_{32}$, $E_{21} = E_{23}E_{31}$. Hence $v_{n}$ commutes with $W$, and $[n] = R(n)v_{n}$.

\emph{Step 5.} If $z \in A$ commutes with $W$ and $Pz = 0$ for a rank-one idempotent $P$ of $W$, then $z = 0$: $1 = \sum_{k}X_{k}PY_{k}$ with $X_{k}, Y_{k} \in W$, so $z = \sum_{k}X_{k}PzY_{k} = 0$. Let $a, b$ be orthogonal nonzero rational bivectors and $c = ab$. From $[a][b] = [c]$ and $R(a)R(b) = R(c)$, $v_{c} = v_{a}v_{b}$, and the relation $\kappa_{a,b} = 0$ of $A$ reads $1 + R(a)v_{a} + R(b)v_{b} + R(c)v_{a}v_{b} = 0$. Multiplying by $P_{a}$, $P_{b}$, $P_{c}$ (on which $R(a), R(b), R(c)$ act by the signs $(+,-,-)$, $(-,+,-)$, $(-,-,+)$) and using the first sentence,
\[
1 + v_{a} - v_{b} - v_{a}v_{b} = 0, \qquad 1 - v_{a} + v_{b} - v_{a}v_{b} = 0, \qquad 1 - v_{a} - v_{b} + v_{a}v_{b} = 0 .
\]
The sum of the first two gives $v_{a}v_{b} = 1$, then $v_{a} = v_{b}$, and the third gives $v_{a} = 1$. Every rational bivector lies in such a triple, so $[n] = R(n) \in W$ for every $n$.

\emph{Step 6.} Let $x = x_{0} + u$ with $u$ a bivector and choose a nonzero rational bivector $m$ orthogonal to $u$; then $xm = x_{0}m + um$ is a bivector and $x = (xm)m^{-1}$. So $G$ is generated by the classes of bivectors, $G_{W} = G$ and $A = W$: $[x] = R(x)$ for all $x$.

\emph{Conclusion of (iii).} A function $f$ with frame sums $0$ is $f(x) = \Lambda(R(x))$ with $\Lambda = f \circ (R|_{W})^{-1}$ linear on $\Mat_{3}(\Q)$. The entries of $R(x_{r})$ are the components of $x_{r}B_{l}\tilde{x}_{r}$ divided by $r \cdot r$, quadratic forms in $r$ over $r \cdot r$; so $f(r) = \sum_{i,j}\tau_{ij}r_{i}r_{j}/(r \cdot r)$ with $\tau_{ij} = \tau_{ji} \in V$, and adding the constant back, $\rho_{ij} = \tau_{ij} + \tfrac14\mu(1)\delta_{ij}$. The ten functions $r_{i}r_{j}/(r \cdot r)$, $i \le j$, are linearly independent on rays (evaluate at the basis vectors and at their pairwise sums), which gives uniqueness of $\rho$. The map $L(x) = \sum_{i,j}\rho_{ij}m(x)_{ij}$ is linear, $L(\Pi_{r}) = \sum\rho_{ij}r_{i}r_{j}/(r \cdot r)$, $L(x^{\dagger}) = L(x)$ because $\rho$ is symmetric, and $L(E) = \mu(E)$ for every projector by (i). Symmetric rational matrices are spanned by the $rr^{T}$, so a linear extension is determined on self-adjoint elements, and the condition $L(x^{\dagger}) = L(x)$ fixes it on the rest.

\emph{(iv).} Let $V = \Q$, $\mu \ge 0$ on atoms and $\mu(1) = 1$. Then $r^{T}\rho r \ge 0$ for every rational $r$, and for any $x$, $L(xx^{\dagger}) = \sum_{i,j}\rho_{ij}(m(x)m(x)^{T})_{ij} = \sum_{k}c_{k}^{T}\rho c_{k} \ge 0$ with $c_{k}$ the columns of $m(x)$; $L(1) = \sum_{i}\rho_{ii} = \mu(1) = 1$. So $L$ is a state. Conversely a state is non-negative on projectors ($E = EE^{\dagger}$), normalised and additive. For the functional $\lambda$: it is linear, so its restriction is an additive weight; the check computes $\lambda(s_{i}s_{j}^{\dagger})$ and $\lambda(\Pi_{(1,0,0,1)}) = -\tfrac14$.

\emph{(v).} Let $r$ be a primitive integer vector with $r \cdot r = d^{2}$, and $t \ge 1$ the number of its odd coordinates. Squares are $0$, $1$ or $4$ modulo $8$, so $r \cdot r \equiv t$ or $t + 4$ modulo $8$, and this is a square modulo $8$ only for $t = 1$ and $t = 4$. Two vectors with one odd coordinate each in the same place, or one of each kind, have an odd inner product, so they are not orthogonal. Hence a frame of such rays consists of four vectors with all coordinates odd, or of four vectors with one odd coordinate each, in four different places, and in both cases $g$ sums to $1$. If $g$ were $\sum\rho_{ij}r_{i}r_{j}/(r \cdot r)$ on these rays, then with $\alpha_{0} = 1$ the four rays $(1, \pm 2, \pm 2, 0)$, each of value $1$, would give $4 = \tfrac49(\rho_{00} + 4\rho_{11} + 4\rho_{22})$, while $\rho_{00} = g(e_{0}) = 1$ and $\rho_{11} = \rho_{22} = 0$ give $\tfrac49$. For other $\alpha \ne (\tfrac14, \tfrac14, \tfrac14, \tfrac14)$ the same four rays around a coordinate $i$ with $\alpha_{i} \ne \tfrac14$ and two other coordinates $j, k$ give $4\alpha_{i} = \tfrac49(\alpha_{i} + 4\alpha_{j} + 4\alpha_{k})$ for each choice of $j, k$, which forces $\alpha_{j} + \alpha_{k} = 2\alpha_{i}$ for all pairs, hence all $\alpha$ equal.
